\section{Proposed Machine Vision-Based Virtual Encoder System}
\label{proposed_encoder}
The proposed virtual encoder operates by mapping sub-pixel visual coordinate transformations into physical metric positions. The core pipeline consists of raw image feature extraction, sub-pixel corner optimization, localized enhancement for homography alignment, pinhole model restoration, and iterative pose estimation.

\subsection{Sub-Pixel Feature Extraction Pipeline}
To preserve native intensity gradient sharpness ($\nabla I$) and avoid spatial boundary shifts, the sub-pixel feature extraction pipeline operates directly on the unenhanced raw grayscale image ($I_{\text{gray}}$). Applying contrast amplification or smoothing filter passes at this initial stage introduces spatial bias and blur artifacts that degrade sub-pixel corner accuracy.

First, the Edge Drawing Lines (EDLines) detector traces continuous pixel chains representing the structural border segments of the chessboard marker on $I_{\text{gray}}$. The line screening process evaluates edge junction topologies to filter background clutter, identifying line clusters intersecting at 3-way or 4-way junctions based on line lengths ($l$), boundary intersection radii ($r$), and angular deviation thresholds ($\theta$). Coarse pixel-level coordinates are initially generated using a localized least-squares line-fitting operator:

\begin{equation}
p = \min_{q_i \in N(p)} \sum \left( g_{q_i}^{\top} (q_i - p) \right)^2
\end{equation}

where $g_{q_i}$ represents the gray gradient vector at an adjacent neighborhood pixel $q_i$ centered around the corner estimation index $p$. Sub-pixel refinement is mathematically resolved by minimizing the projection error where the vector connecting the true corner to neighborhood structures is perfectly orthogonal to the local gradient, expressed analytically as:

\begin{equation}
p = \left(\sum_{q_i \in N(p)} g_{q_i} g_{q_i}^{\mathsf{T}}\right)^{-1} \cdot \sum_{q_i \in N(p)} \left(g_{q_i} g_{q_i}^{\mathsf{T}}\right) q_i
\end{equation}

%%%%%%%


\subsection{Homography Alignment and Corner Sorting}
To arrange the 49 extracted sub-pixel corners into a unified coordinate grid, three asymmetric internal circular elements located at the marker's center serve as spatial anchors. Unlike the sub-pixel optimization phase, anchor detection requires high local contrast and noise suppression to prevent false circle candidates. Therefore, image enhancement is selectively applied strictly within this sorting phase.

The grayscale image is first enhanced using Contrast Limited Adaptive Histogram Equalization (CLAHE) to amplify local contrast boundaries around the central anchor patterns. To suppress high-frequency noise and amplification artifacts induced by CLAHE, a 2D Gaussian blur convolves the enhanced image using a symmetric $5 \times 5$ kernel:

\begin{equation}
G_2(x,y) = \frac{1}{2\pi\sigma^2} e^{-\frac{x^2+y^2}{2\sigma^2}}
\end{equation}

where $\sigma$ acts as the standard deviation controlling the smoothing threshold. Isolating CLAHE and Gaussian filtering to this alignment stage enables robust anchor extraction via the Hough Gradient Circle Transform (or EDCircles) without corrupting the native pixel gradients used in sub-pixel corner optimization.

Following anchor centroid isolation, candidate corners are filtered to select the central $7 \times 7$ grid points. A homography transformation matrix $H$ is calculated using the four extreme boundary corners to project perspective coordinates into a structured $7 \times 7$ planar layout, effectively neutralizing out-of-plane rotation drifts and ordering corners in row-major sequence.

\subsection{Projection Mapping and Iterative SolvePnP}
The physical tracking domain relies on standard pinhole camera geometry. A 3D spatial coordinate vector $[X, Y, Z]^T$ defined in the world frame is projected onto the 2D image plane pixel index $[u, v, 1]^T$ via the intrinsic matrix ($K$) and extrinsic transformation profile ($[R|T]$):

\begin{equation}
s \begin{bmatrix} u \\ v \\ 1 \end{bmatrix} = \begin{bmatrix} f_x & 0 & c_x \\ 0 & f_y & c_y \\ 0 & 0 & 1 \end{bmatrix} \begin{bmatrix} R_{11} & R_{12} & R_{13} & T_x \\ R_{21} & R_{22} & R_{23} & T_y \\ R_{31} & R_{32} & R_{33} & T_z \end{bmatrix} \begin{bmatrix} X \\ Y \\ Z \\ 1 \end{bmatrix}
\end{equation}

where $s$ scales the perspective mapping, $f_x, f_y$ represent directional focal lengths, and $c_x, c_y$ pinpoint the sensor's principal optical center. Lens distortion is mitigated using an all-encompassing radial and tangential mapping law:

\begin{equation}
\begin{cases} 
x' = x(1 + k_1r^2 + k_2r^4) + 2p_1xy + p_2(r^2 + 2x^2) \\ 
y' = y(1 + k_1r^2 + k_2r^4) + p_1(r^2 + 2y^2) + 2p_2xy 
\end{cases}
\end{equation}

with $r^2 = x^2 + y^2$, mapping structural sensor anomalies through coefficients $k_n$ and $p_n$. 

The actual translation vector $T = [T_x, T_y, T_z]^T$ provides direct millimeter-scale metric tracking data, eliminating external scale factors. Finding the optimal camera pose is framed as a non-linear optimization problem solved using the iterative Levenberg-Marquardt (LM) method. The update vector ($\Delta x$) is updated at each step by adjusting a trust-region damping parameter ($u$):

\begin{equation}
\Delta x = (G^T G + uI)^{-1} G^T f(x)
\end{equation}

where $G$ denotes the calculated Jacobian matrix of the projection functions and $I$ defines a standard identity matrix.